\section{Заключение}

В ходе прохождения учебной ознакомительной практики выполнены
и протестированы 15 индивидуальных задач (вариант~4), охватывающих
все темы методических указаний: линейные и разветвляющиеся
алгоритмы, циклические и итерационные вычисления, операции над
массивами, векторами и матрицами, линейный поиск, арифметика,
геометрия и теория множеств, линейная алгебра, обработка
символьной информации, аналитическая геометрия, кривые и
поверхности второго порядка.

\medskip
\noindent\textbf{Для каждой задачи выполнено:}
\begin{itemize}
  \item приведено математическое обоснование решения;
  \item построена блок-схема алгоритма средствами~TikZ;
  \item разработан программный код на языке~C с пояснительными
        комментариями;
  \item программа протестирована на нескольких наборах входных
        данных, включая граничные случаи.
\end{itemize}

\medskip
\noindent\textbf{Практическая значимость.} Выполненные задачи
демонстрируют владение базовыми и продвинутыми приёмами
программирования на языке~C: работу с условными конструкциями
и циклами, массивами произвольной размерности, рекурсией
(вычисление определителя минора в теме~9), битовыми масками
для перебора подмножеств (тема~8), а также реализацию
геометрических и алгебраических вычислений с использованием
библиотеки~\texttt{math.h}.

\medskip
\noindent\textbf{Выводы.} Все программы скомпилированы без
предупреждений компилятора (флаги \texttt{-Wall~-Wextra}) и дают
корректные результаты на контрольных примерах, что подтверждает
правильность реализованных алгоритмов. Практика позволила закрепить
навыки структурного программирования, формального описания
алгоритмов средствами блок-схем и оформления технической
документации в системе~\LaTeX{}.
